\addchap{ABSTRACT}

Methane is one of the most powerfull greenhouse gases contributing to climate change with about 25 times as much warming potential as carbondioxide over the span of one hundred years. To reduce the methane content in the atmosphere, it is highly important to understand methane sources and sinks. While a lot of methane is emitted by industries and waste sites, there are also natural emissions from, for example, wetlands. This shows the diverse environments in which monitoring is necessary. For this reason, current-day large-scale monitoring is done from space, relying on atmospheric models to calculate the methane concentration profile from the averaged column density. These atmospheric models require in-situ measurement data to make them more accurate. However, not many missions have done so for methane in the stratosphere, which makes the models less reliable overall and therefore renders monitoring less reliable as well. \\
Furthermore, commonly used laser-based technologies measuring methane with high precision are often very expensive and complex. This prevents more in-situ measurements from being recorded in the stratosphere.\\  
Methane InfraRed Absorption Gas Experiment (MIRAGE) aims to demonstrate a low-cost, low-complexity alternative to laser-based technologies to measure stratospheric methane concentrations, using an off-the-shelf Non-Dispersive InfraRed (NDIR) sensor core. To increase the accuracy, temperature-controlled ambient air is pressurised before feeding it into the optical chamber of the NDIR sensor. With this experiment, NDIR technology will be tested in regards to its limitations and future potential for use in atmospheric methane monitoring.
%This is done by pressurizing temperature-controlled ambient air before feeding it into the optical chamber of the NDIR sensor.
%With this setup it is expected to achieve an absolute concentration accuracy of \dos{$\mathrm{\Delta}c \leq 0.1$ ppm}.


%\dos{EXPAND THIS. Expansion is not a required action, but it was commented that our abstract is very short.}